\documentclass[11pt,a4paper]{article}
\usepackage[margin=1in]{geometry}
\usepackage{hyperref}
\usepackage{enumitem}
\usepackage{graphicx}
\usepackage{xcolor}

% -----------------------------------------------------
% bvraghav-tiet-ucs503p-article.sty
% -----------------------------------------------------

% Variable: Lab Instructor
\makeatletter \newcommand{\BvrSetLabInstructor}[1]{%
  \gdef\Bvr@LabInstructor{#1}}
% default
\newcommand{\Bvr@LabInstructor}{Paramveer Sidhu}
% public accessor
\newcommand{\BvrLabInstructorValue}{\Bvr@LabInstructor}
\makeatother

% Add subtitle
\usepackage{titling}
% -- Set title bold
\pretitle{\begin{center}\bfseries\fontsize{18bp}{18bp}\selectfont}
\makeatletter
\newcommand{\BvrSubtitle}[1]{%
  \posttitle{%
    \par\end{center}
    \begin{center}\large#1\end{center}
    \vskip 0.5em}%
}
\makeatother

% Hints in the section title(s)
\newcommand{\BvrHint}[1]{%
  {\normalfont\normalsize\itshape\hfill #1}}

% -----------------------------------------------------

\title{MindPulse: A Privacy-First Multimodal Cognitive Fatigue and Burnout Monitoring Pipeline}

\BvrSetLabInstructor{Paramveer Sidhu}

\BvrSubtitle{%
  Project Proposal for UCS503P  \\
  Submitted to: \BvrLabInstructorValue \\ \bfseries
  Thapar Institute of Engineering and Technology% 
}

\author{
  Aditi Sinha, 3C11%
  \thanks{Roll No: \texttt{1024030024}, %
    Email: \texttt{asinha\_be24\@thapar.edu}}
  \and 
  Sajneet Kaur Grewal, 3C11%
  \thanks{Roll No: \texttt{1024030016}, %
    Email: \texttt{sgrewal\_be24\@thapar.edu}}
}

\date{\today}

\begin{document}
\maketitle

\section{Higher-order goal%
  \BvrHint{\ldots secondary framing}}
This project aims to enhance student and developer wellness by reducing cognitive exhaustion and screen-induced burnout. While workplace fatigue is a widespread behavioral issue, the immediate engineering objective is to translate \emph{passive fatigue detection and recovery scheduling} into a measurable, lightweight, and deployable software pipeline.

\section{Time-to-value %
\BvrHint{\ldots primary heuristic}}
\begin{itemize}[leftmargin=0.7cm]
    \item We will deliver meaningful increments quickly by prototyping the core session-logging and fatigue-scoring pipeline within a 2-week iteration window.
    \item We will prioritize fast feedback over custom model research by using pre-trained feature extractors and weekly CI-driven automated test releases.
    \item Success is defined by what can be scored, logged, and visualized in the first iteration, and incrementally expanded with multimodal inputs thereafter.
\end{itemize}

\section{Problem Statement}
University students and knowledge workers experience severe productivity drops and digital fatigue during prolonged, uninterrupted screen sessions. Common pain points include:
\begin{itemize}[leftmargin=0.7cm]
    \item \textbf{Delayed self-awareness:} users rarely recognize cognitive saturation until headaches, attention lapses, or coding blunders occur.
    \item \textbf{Intrusive monitoring:} existing surveillance tools record invasive video feeds or keystrokes, raising severe privacy concerns.
    \item \textbf{Lack of actionable recovery loops:} traditional wellness tools provide static reminders rather than contextual break recommendations based on session strain.
    \item \textbf{No verifiable audit of session load:} users cannot correlate long study/work sprints with declining focus over time.
\end{itemize}

\textbf{Why this matters now (contextual relevance):} In remote learning and intensive engineering environments, a non-invasive, privacy-preserving fatigue sentinel prevents burnout and sustains long-term academic productivity.

\section{Proposed Solution %
\BvrHint{\ldots Software Proposition}}
\subsection{Overview}
Build a \textbf{web-based} ``Cognitive Fatigue Sentinel'' platform with the following components:
\begin{itemize}[leftmargin=0.7cm]
    \item \textbf{Client-side telemetry collector:} processes lightweight, on-device metrics (eye blink rate/EAR via webcam and typing cadence rhythm) without uploading raw media to servers.
    \item \textbf{Self-reported strain pulse:} quick 2-click micro-survey capturing subjective focus level.
    \item \textbf{Deterministic composite scoring engine:} combines real-time session duration, micro-survey inputs, and edge blink metrics into an actionable Fatigue Index ($0$--$100$).
    \item \textbf{Live wellness dashboard:} displays current strain tier (Optimal, Moderate, Fatigued) and smart recovery countdowns.
    \item \textbf{Smart intervention trigger:} automated break prompts and micro-recovery exercises dispatched upon threshold violations.
\end{itemize}

\subsection{Core workflow}
\begin{enumerate}[leftmargin=0.7cm]
    \item User starts a work/study session via the web dashboard.
    \item Browser extracts local blink rate telemetry via client-side libraries and logs session interval metrics.
    \item Ingestion API consumes structured numerical telemetry payloads (no raw video or text stored).
    \item Rule engine computes the composite Fatigue Index and flags threshold breaches.
    \item Dashboard alerts user when fatigue levels exceed safe thresholds and recommends specific recovery intervals.
\end{enumerate}

\subsection{Operational constraints for an educational, campus setting}
\begin{itemize}[leftmargin=0.7cm]
    \item Ensure zero raw video or keystroke content leaves the user's local browser to preserve complete privacy.
    \item Avoid building deep neural networks from scratch; rely on mature client-side vision primitives and deterministic scoring heuristics.
    \item Design for low CPU utilization so background monitoring does not degrade system performance.
\end{itemize}

\section{Solution Approach %
\BvrHint{\ldots Engineering focus}}
This is an engineering course project, so we focus on data pipelines, client-side efficiency, and automated testing rather than novel algorithmic research.

\subsection{Fatigue scoring heuristics %
\BvrHint{\ldots non-research}}
Fatigue evaluation will be deterministic and rule-weighted:
\begin{itemize}[leftmargin=0.7cm]
    \item Use pre-trained client-side landmark extractors (e.g., MediaPipe FaceMesh) to derive Eye Aspect Ratio (EAR) locally.
    \item Calculate weighted score vectors across session duration, blink frequency drop, and self-reported focus markers.
    \item Maintain deterministic boundary thresholds (e.g., Score $> 75 \rightarrow$ High Fatigue) for predictable testing.
\end{itemize}

\subsection{Web architecture %
\BvrHint{\ldots web-based solution}}
A standard 3-tier web architecture:
\begin{itemize}[leftmargin=0.7cm]
    \item \textbf{Frontend:} responsive single-page web app with client-side vision extraction and live charting.
    \item \textbf{Backend API:} REST endpoints for ingesting numerical telemetry records, user sessions, and threshold alerts.
    \item \textbf{Database:} relational database (SQLite/PostgreSQL) storing session timestamps, fatigue scores, and recovery logs.
\end{itemize}

\subsection{CI/CD and fast delivery enabler}
CI/CD will be used to:
\begin{itemize}[leftmargin=0.7cm]
    \item Automatically run unit test suites on score-weighting algorithms and schema validation on every git push.
    \item Verify that invalid or negative telemetry values are rejected cleanly.
    \item Deploy verified builds to cloud staging environments automatically.
\end{itemize}

\section{Evaluation Criterion %
\BvrHint{\ldots measurable and attributable}}
We define success using measurable metrics tied to the core workflow.

\subsection{Primary evaluation metric}
\textbf{Intervention Trigger Latency (ITL):} time between telemetry ingestion crossing a fatigue threshold and the dispatch of a recovery notification.
\begin{itemize}[leftmargin=0.7cm]
    \item Target: median ITL $\le$ 1.5 seconds during active monitoring sessions.
    \item Attribution: measured from payload receipt timestamp to alert dispatch log in the database.
\end{itemize}

\subsection{Secondary evaluation metrics}
\begin{itemize}[leftmargin=0.7cm]
    \item \textbf{Deterministic scoring accuracy:} 100\% pass rate on edge-case test suites (verifying correct state transitions at boundary scores 50, 75, 90).
    \item \textbf{Data efficiency:} telemetry payload sizes restricted to $< 2\text{ KB}$ per ping to maintain lightweight operation.
    \item \textbf{Client execution overhead:} client-side feature extraction consuming $< 5\%$ average CPU on standard laptops.
    \item \textbf{Pipeline reliability:} $\ge$ 99\% CI test pass rate prior to release.
\end{itemize}

\subsection{Pilot validation plan %
\BvrHint{\ldots high-level}}
\begin{itemize}[leftmargin=0.7cm]
    \item Run automated mock telemetry scripts simulating 2-hour continuous study sessions with escalating fatigue patterns.
    \item Conduct user testing with 10 student volunteers across 1-hour study blocks to evaluate alert timeliness and UI clarity.
\end{itemize}

\section{Scalability %
\BvrHint{\ldots with foresight}}
The architecture is designed to scale across multiple users and institutions.

\begin{enumerate}[leftmargin=0.7cm]
    \item \textbf{Scalable (theoretically):} The system handles high concurrent user loads by:
    \begin{itemize}[leftmargin=0.7cm]
        \item delegating all visual/computational extraction to client browsers,
        \item keeping backend APIs stateless and lightweight (processing only numerical scores),
        \item indexing database records on session IDs and timestamps for instant history querying.
    \end{itemize}
    
    \item \textbf{Operationally deployable (practically):} The system runs reliably on common infrastructure:
    \begin{itemize}[leftmargin=0.7cm]
        \item managed relational database,
        \item lightweight containerized deployment (Docker),
        \item automated GitHub Actions pipeline for continuous integration and testing.
    \end{itemize}
\end{enumerate}

\section{Engine availability heuristic}
A mature set of stable libraries is available to implement this without research delays:
\begin{itemize}[leftmargin=0.7cm]
    \item Client-side vision primitives (MediaPipe / FaceMesh via JavaScript).
    \item Backend REST API framework (FastAPI or Express.js).
    \item Relational database engine (SQLite for dev, PostgreSQL for staging).
    \item Dashboard charting engines (Chart.js / Recharts).
    \item Test runner and CI engine (Pytest / Jest and GitHub Actions).
\end{itemize}
These mature tools enable the team to focus on input reliability, privacy guarantees, and UX rather than custom machine learning model training.

\section{Project scope and deliverables %
\BvrHint{\ldots iteration-friendly}}
\subsection{Initial deliverable %
\BvrHint{\ldots first iteration}}
\begin{itemize}[leftmargin=0.7cm]
    \item Telemetry ingestion endpoint (\texttt{POST /api/fatigue}) with strict schema validation
    \item Deterministic scoring engine computing fatigue index from basic inputs (session duration + micro-survey)
    \item Database schema storing user session records and calculated scores
    \item Web dashboard displaying active session status with a Safe/Fatigued visual indicator
    \item CI pipeline: automated linting and unit tests verifying score calculation logic
\end{itemize}

\subsection{Subsequent deliverables}
\begin{itemize}[leftmargin=0.7cm]
    \item Client-side eye aspect ratio (EAR) integration via MediaPipe for passive blink tracking
    \item Historical session analytics showing focus and fatigue trends over time
    \item Interactive break scheduler with micro-recovery guidance prompts
    \item Automated synthetic data generator for stress-testing pipeline edge cases
    \item CSV/PDF session report export for personal habit tracking
\end{itemize}

\section{Risks and mitigations}
\begin{itemize}[leftmargin=0.7cm]
    \item \textbf{User privacy concerns:} guarantee that no video frames, images, or raw text are transmitted or stored; only derived numerical metrics are sent to the backend.
    \item \textbf{Camera lighting sensitivity:} allow the system to gracefully fall back on session duration and survey inputs if optical conditions are poor.
    \item \textbf{Alert fatigue:} introduce cooldown timers to prevent repeated notifications within short time windows.
    \item \textbf{Evaluation measurement ambiguity:} use automated server timestamps to strictly measure alert response latencies.
\end{itemize}

\section{Summary}
This project proposes MindPulse, a privacy-preserving web platform designed to detect cognitive fatigue and recommend timely recovery breaks during extended computer work. It meets the course criteria by delivering a functional 2-week MVP, using CI automation to guarantee calculation correctness, and leveraging mature client-side engines to bypass model training bottlenecks. It focuses on engineering reliability, data privacy, and measurable intervention latency.

\end{document}
